\section{Related Work}

\paragraph{Knowledge Distillation.}
Knowledge distillation (KD)~\citep{hinton2015distilling} transfers knowledge from a large model to a smaller one by training a student network on the soft output distributions of a teacher.
For autoregressive sequence models,~\citet{kim2016sequence} extended this to sequence-level distillation by training students on teacher-generated outputs, establishing the dominant off-policy distillation baseline~\citep{sanh2019distilbert,jiao2020tinybert,wang2020minilm}.
In parallel, supervised fine-tuning (SFT) has been directly applied to improve performance on a variety of downstream tasks~\citep{JMLR:v25:23-0870,sanh2021multitask,wei2021finetuned}.
A fundamental limitation shared by all off-policy approaches is the train-inference distribution mismatch. The student is optimized on teacher-generated or reference sequences, but must generate from its own distribution at inference, which is an instance of the exposure bias~\citep{bengio2015scheduled} that accumulates errors over long generations.
This mismatch motivates shifting distillation to the student's own on-policy distribution, which is the central idea behind on-policy distillation.


\paragraph{On-Policy Distillation.}
MiniLLM~\citep{gu2023minillm} first formalized on-policy distillation (OPD) for LLMs under a reverse KL objective optimized via policy gradient, arguing that reverse KL's mode-seeking behavior prevents the student from spreading probability mass over regions the teacher considers unlikely.
GKD~\citep{agarwal2024policy} introduced a unified framework interpolating between on-policy and off-policy data across multiple divergences, demonstrating consistent gains over other KD baselines.
\citet{yang2026learning} later formalized OPD theoretically as a special case of dense KL-constrained RL, showing that the teacher's per-token log-ratio constitutes an implicit reward and that scaling this reward beyond its standard weight can push the student past the teacher's performance boundary.
OPD has since been adopted in industry post-training pipelines~\citep{yang2025qwen3,lu2025onpolicydistillation,zeng2026glm,xiao2026mimo,ko2026scaling,jin2026entropy,jang2026stable, fu2026revisiting,yang2026learning}, and extended to scalable self-distillation~\citep{hubotter2026reinforcement,zhao2026self,he2026far,shenfeld2026self,ye2026policy,sang2026crispcompressedreasoningiterative,kim2026does,ye2026online,yang2026self,li2026unifying,zhao2026selfdistillationmultitokenprediction,ding2026hdpo}, where a single model acts as its own teacher by conditioning on privileged information such as ground-truth solutions or execution feedback.
Despite this growing body of work, existing studies focus on demonstrating OPD's promise, such as dense rewards and mitigated exposure bias, across varied objectives, tasks, and teacher-student pairs, without systematically analyzing when or why OPD fails.


\paragraph{Capacity Gap and Distillability.}
A recurring observation in knowledge distillation is that large teacher-student capacity gaps can degrade or even reverse the benefit of distillation.
\citet{cho2019efficacy} demonstrate that distillation can hurt student performance when the teacher is substantially more capable, and~\citet{mirzadeh2020improved} propose an intermediate-sized teacher assistant to bridge the gap.
\citet{busbridge2025distillation} provide a quantitative treatment via distillation scaling laws, showing that student loss follows a power law as a function of teacher quality, student size, and data volume, identifying a U-shaped capacity regime where teacher over-capability degrades distillation efficiency.
For LLM reasoning,~\citet{li2025small} document a ``learnability gap'' showing that training small models on long chain-of-thought traces from strong reasoning teachers consistently underperforms simpler approaches, suggesting that the reasoning complexity of teacher outputs must be matched to student capacity.
These findings call for caution regarding the universality of distillation. However, the existing analyses have largely centered on off-policy knowledge distillation. In particular, the issues of capacity gap and distillability in OPD remain underexplored.











